\subsection{KL-regularized control and preference optimization}
\label{sec:C08:dpo}
Consider one context $x$ and a response distribution $\pi$. With reference $\pi_0$, reward $r(x,y)$, and $\beta>0$, maximize $\mathbb E_\pi r-\beta\operatorname{KL}(\pi\|\pi_0)$. Assuming the reference has support on feasible responses and the partition function is finite, a Lagrange-multiplier calculation gives
\begin{equation}
 \pi^*(y\mid x)=\frac{\pi_0(y\mid x)e^{r(x,y)/\beta}}
 {Z(x)}.
\end{equation}
Rearranging yields $r(x,y)=\beta\log[\pi^*(y\mid x)/\pi_0(y\mid x)]+\beta\log Z(x)$. Under a Bradley--Terry preference model, reward differences cancel the context-only partition term. This is the algebraic basis of direct preference optimization~\cite{rafailov2023dpo}.

For a preferred completion $y^+$ and dispreferred completion $y^-$, define
\begin{align}
 d_\theta={}&\log\frac{\pi_\theta(y^+\mid x)}{\pi_0(y^+\mid x)}
 -\log\frac{\pi_\theta(y^-\mid x)}{\pi_0(y^-\mid x)},\\
 \mathcal L_{\mathrm{DPO}}={}&-\mathbb E\log\sigma(\beta d_\theta).
\end{align}
The per-example gradient is $-\beta\sigma(-\beta d_\theta)\nabla_\theta d_\theta$. Incorrectly ranked pairs receive stronger weight under this formulation. The derivation does not prove that observed preferences are generated by one stationary scalar reward, nor that the model class can realize the optimum. Cyclic preferences, label noise, and distribution shift remain outside the simple equivalence argument.

Agent preferences can compare trajectories with different external states, which complicates the meaning of a pair. A shorter trace may be preferable only because it skipped an expensive but necessary check. A successful trace can violate a permission boundary that a failed compliant trace respected. A single binary preference can conflate task success, safety, cost, and presentation. The dataset should either encode the tradeoff deliberately or retain separate labels. Hard execution constraints should not be converted into a negotiable preference merely because the training method uses pairwise comparisons.

Support is another practical limit. If the reference assigns negligible probability to a useful action sequence, regularized optimization may struggle to discover it from static preference data. Offline pairs also provide little evidence about actions absent from the dataset. Online environment interaction can expose new failures, but introduces cost and safety requirements. DPO and reinforcement learning therefore optimize different data interfaces even when their idealized objectives are related. Comparing them requires matching not only model size but data exposure, online calls, reward access, and the opportunity to discover new trajectories.

